\documentclass[11pt]{article}

\usepackage[utf8]{inputenc}
\usepackage[T1]{fontenc}
\usepackage{lmodern}
\usepackage{amsmath,amssymb}
\usepackage{booktabs}
\usepackage{graphicx}
\usepackage{hyperref}
\usepackage{url}
\usepackage{microtype}
\usepackage[margin=1in]{geometry}

\hypersetup{
  colorlinks=true,
  linkcolor=blue,
  citecolor=blue,
  urlcolor=blue
}

\title{RWKV-7 State Semantic Embedding: A Unified Supervised Projection
Framework for Three Tasks, with a Production Agentic Router}

\author{Zhenyuan Gu\thanks{Corresponding author. Email:
\url{admin@ai00-x.com}. Work done as an independent researcher; production
deployment in the open-source ai00-x client
(\url{https://github.com/cgisky1980/ai00-x-client}).}}

\date{August 21, 2026}

\begin{document}

\maketitle

\begin{abstract}
RWKV-7, as a modern linear RNN architecture, contains two levels of internal
states: (1) the \emph{WKV state}---a recurrent memory matrix shaped
$[\text{num\_heads}, \text{key\_dim}, \text{value\_dim}]$ that is the core of
the Delta Rule; and (2) the \emph{hidden state}---the output vector of each
TMix layer. This paper finds that the hidden state extracted by albatross
(BlinkDL's inference engine) contains rich semantic information but exhibits
severe anisotropy, preventing unsupervised methods (KMeans, cosine similarity)
from effectively mining it---achieving only $v$-measure $=0.44$ on 20
Newsgroups (sklearn full-text version) clustering and Spearman $=0.46$ on
STS-B semantic similarity. As a comparison, we also tried various aggregation
methods on the WKV state (Q-Readout, row\_sum, diag, etc.), yielding even
worse clustering performance (best 0.11), indicating that the hidden state is
more suitable for semantic extraction than the WKV state. This paper presents
a key insight: \textbf{hidden state contains semantic information
(supervised MLP classification achieves 0.93), but unsupervised methods
cannot extract it; supervised data is needed to train task-specific
projectors to unlock its semantic potential.} Based on this insight, all
three tasks achieve the following results: \textbf{(1) Semantic Similarity}:
trained an STS-specific projector using 46.9k labeled pairs from
NLI/STS/SICK-R (strictly deduplicated, removing 1,249 leak pairs overlapping
with STS-B dev/test), achieving Spearman $=0.8188$; \textbf{(2) Topic
Clustering}: validated on two versions of 20 Newsgroups---MTEB short-text
version (59,545 samples) supervised projection $v$-measure $=0.9506$ ($+217\%$
over the unsupervised baseline 0.2912); sklearn full-text version (strict
dedup + stratified 70/15/15 split) SupCon loss $v$-measure $=0.6660$ ($+50\%$
over the unsupervised baseline 0.4434); \textbf{(3) Task Classification}:
achieved test accuracy $=0.9325$ on an independent test set using
Hidden+MLP. It should be emphasized that STS and clustering tasks use
supervised projectors and are thus supervised methods; the comparison with
unsupervised baselines is only to validate the claim that hidden state
requires supervised projection, not to claim superiority over unsupervised
methods. Beyond offline benchmarks, the classification pipeline is deployed
as the production smart router of an open-source desktop assistant: on a 2.9B
int8 backbone with Chinese-language augmentation, held-out test accuracy
reaches 0.960 (0.971 end-to-end through the deployed engine) at 77\,ms mean
latency, with zero runtime dependencies beyond the already-resident local
model. All offline results are based on the albatross inference engine
(source code unmodified), a 0.4B RWKV-7 model, and CPU training (parameters
$\sim$3.15M), with all code open-source and reproducible.
\end{abstract}

\section{Introduction}

\subsection{Background}

RWKV-7 \cite{rwkv7} is a modern linear RNN that achieves $O(L)$ complexity
sequence modeling through the Delta Rule \cite{deltarule}. It contains two
levels of internal states:

\begin{itemize}
  \item \textbf{WKV state} $S \in \mathbb{R}^{\text{num\_heads} \times
  \text{key\_dim} \times \text{value\_dim}}$: recurrent memory matrix, the
  core of the Delta Rule, recursively updated at each time step.
  \item \textbf{Hidden state} $h \in \mathbb{R}^{\text{hidden\_dim}}$: the
  output vector of each TMix (time mix) and CMix (channel mix) layer, an
  intermediate representation of forward propagation.
\end{itemize}

This paper studies the semantic extraction capability of the \textbf{hidden
state}. The WKV state was also experimented with as a baseline
(\S\ref{sec:unsupervised-comparison}), but its clustering performance was
much lower than the hidden state (0.11 vs 0.47). Albatross \cite{albatross}
is an efficient inference engine developed by BlinkDL, providing
CUDA-accelerated WKV kernels and batch concurrent inference capability.

\subsection{Problem}

Although the albatross-extracted hidden state contains semantic information,
it suffers from severe anisotropy:

\begin{table}[h]
\centering
\begin{tabular}{lll}
\toprule
Task & Unsupervised Method & Metric \\
\midrule
Topic Clustering & Hidden + standardize + KMeans & $v$-measure = 0.47 (sklearn full-text) \\
Semantic Similarity & Hidden cosine & Spearman = 0.46 \\
Task Classification & Hidden + MLP (supervised) & test acc = 0.93 \\
\bottomrule
\end{tabular}
\caption{Anisotropy of the raw hidden state across tasks.}
\end{table}

The supervised MLP classification achieving 0.93 proves that the hidden state
does contain semantic information; however, unsupervised methods perform
poorly on clustering and similarity tasks, indicating that \textbf{information
requires nonlinear transformation to be extracted}.

\subsection{Contributions}

\begin{enumerate}
  \item \textbf{Key Insight}: the anisotropy of the albatross hidden state
  can be mitigated through supervised projectors without modifying the
  inference engine.
  \item \textbf{Task-Specific Projection}: STS learns similarity ranking
  (46.9k pairs, strictly deduplicated); clustering learns class separation
  (12.8k labeled samples, sklearn full-text version with strict
  deduplication + split)---the two objectives differ and cannot be mixed.
  \item \textbf{Strict Experimental Paradigm}: STS training data is strictly
  deduplicated against STS-B dev/test (removing 1,249 overlapping pairs);
  clustering uses sklearn full-text 20NG with strict deduplication +
  stratified 70/15/15 split, where train is for training, dev for
  best\_state selection, and test for held-out evaluation; classification
  adds an independent test set, with head selection using only dev.
  \item \textbf{Three Tasks Results}: STS Spearman $=0.8188$; clustering
  validated on two datasets (MTEB short-text $v=0.9506$, sklearn full-text
  SupCon $v=0.6660$, $+217\%$/$+50\%$ over unsupervised baselines
  respectively); classification test acc $=0.9325$. Note that STS and
  clustering use supervised projectors and are thus supervised methods.
  \item \textbf{Systematic Unsupervised Comparison}: 7 categories and 30+
  unsupervised methods (KMeans/PCA/UMAP/Whitening/DeepCluster) peak at only
  0.33 on the MTEB short-text version, empirically confirming the limitation
  of unsupervised methods and supporting the necessity of supervised
  projection.
  \item \textbf{Pure Python Implementation}: based on the albatross inference
  engine, bucketed by length for concurrency, 250 samples/s.
  \item \textbf{Production Validation}: the classification head ships as the
  smart router of an open-source desktop assistant---on a 2.9B int8 backbone
  it reaches 0.960 test accuracy and 77\,ms mean latency with zero extra
  runtime dependencies; two deployment findings (next-token logits of an
  untuned base model are unusable for tier classification; a
  language-distribution pitfall in training data causes a
  ``Chinese\,$\to$\,R0'' shortcut) are documented for practitioners.
\end{enumerate}

\section{Related Work}

\subsection{RWKV and Albatross}

RWKV \cite{rwkv7} combines the Transformer's parallel training with the
RNN's inference efficiency. RWKV-7 introduces the Delta Rule update:
\begin{equation}
S_t = \operatorname{Diag}(w_t) S_{t-1} + \beta_t (v_t -
S_{t-1}^{\top} k_t) k_t^{\top}.
\end{equation}
Albatross \cite{albatross} is the official PyTorch inference implementation,
providing CUDA WKV kernels and \texttt{forward\_seq\_batch} for batch
concurrent inference.

\subsection{Embedding Evaluation Benchmarks}

\begin{itemize}
  \item \textbf{MTEB} \cite{mteb}: Massive Text Embedding Benchmark.
  \item \textbf{STS-Benchmark} \cite{stsb}: Semantic Textual Similarity,
  Spearman correlation.
  \item \textbf{20 Newsgroups} \cite{20ng}: 20-category topic document
  clustering.
  \item \textbf{NLI / SICK-R} \cite{sick}: Natural Language Inference and
  sentence pair scoring.
\end{itemize}

\subsection{Anisotropy and Contrastive Learning}

\begin{itemize}
  \item \textbf{AnglE} \cite{angle}: angle-based contrastive loss,
  mitigating embedding anisotropy.
  \item \textbf{SimCSE} \cite{simcse}: contrastive learning to alleviate
  anisotropy, but requires many negative samples.
  \item \textbf{Whitening} \cite{whitening}: linear transformation to
  eliminate anisotropy, but destroys nonlinear structure.
\end{itemize}

\subsection{Agentic Model Routing}

SquillaRouter \cite{squillarouter} routes each agent turn across four cost
tiers (C0--C3) using 390 hand-crafted features with a LightGBM ranker and
ONNX inference; its technical report shows 43--90\% realized-cost reduction
on agentic benchmarks while preserving task quality. Our production
deployment (\S\ref{sec:production}) adopts the same four-tier philosophy but
replaces hand-crafted features and external runtimes with the native hidden
state of the already-resident local model.

\section{Method}

\subsection{Albatross Feature Extraction}

\textbf{Model}: RWKV-7-g1 0.4B (hidden=1024, 24 layers, 16 heads,
head\_size=64).

\textbf{Features}:
\begin{itemize}
  \item \textbf{Hidden state} $h \in \mathbb{R}^{1024}$: mean pooling of the
  last layer's FFN output.
  \item \textbf{WKV state} $S \in \mathbb{R}^{16 \times 64 \times 64}$:
  recurrent memory matrix at layer 12.
\end{itemize}

\textbf{Batch Concurrent Extraction}: bucketed by token length, sequences
within each bucket have identical length (no padding needed), using
albatross's \texttt{forward\_seq\_batch} for concurrent inference at 250
samples/s (vs 70 samples/s for a single sequence, 3.6$\times$ speedup).

\textbf{Important Note}: this work entirely uses the albatross inference
engine without modifying any source code. The feature extraction script only
replicates the \texttt{forward\_seq\_batch} loop structure to simultaneously
extract intermediate-layer state and hidden (the original function only
returns logits).

\subsection{Unified Projection Framework}

Core idea: use a 2-hidden-layer MLP to project the hidden state into a
128-dimensional semantic space:
\begin{equation}
\text{emb} = \operatorname{normalize}(\operatorname{MLP}(\operatorname{BatchNorm}(h))).
\end{equation}

\textbf{MLP Architecture} (parameters $\sim$0.86M):
\begin{verbatim}
Input(1024) -> BatchNorm -> Linear(1024->512) -> GELU -> LayerNorm
           -> Dropout(0.2) -> Linear(512->512) -> GELU -> LayerNorm
           -> Dropout(0.2) -> Linear(512->128) -> L2 Normalize
\end{verbatim}

\textbf{Training objectives differ by task} (key design):

\begin{table}[h]
\centering
\begin{tabular}{llll}
\toprule
Task & Training Data & Loss & Temperature $\tau$ \\
\midrule
STS & 46.9k labeled pairs (scores 0--5, deduplicated) & AnglE (regression) & 0.50 \\
Clustering & 12.8k labeled samples (score 0/1, sklearn split) & AnglE (contrastive) & 0.30 \\
\bottomrule
\end{tabular}
\caption{Task-specific training objectives.}
\end{table}

\textbf{Why not mix}: STS learns ``similarity ranking'' (relative distance),
while clustering requires ``class separation'' (absolute distance).
Experiments show that transferring the STS projector to clustering fails
($v$-measure drops from 0.34 to 0.14, see \S\ref{sec:ablation}).

\subsection{Task 1: Semantic Similarity (STS)}

\textbf{Training Data} (total 46,898 pairs after strict deduplication).
Originally 48,147 pairs; after removing pairs overlapping with STS-B
dev/test, 46,898 pairs remain (1,249 pairs removed, 2.59\%):

\begin{table}[h]
\centering
\begin{tabular}{lrrll}
\toprule
Dataset & Original & After dedup & Score range & Source \\
\midrule
STS-B train & 5,749 & 5,076 & 0--5 & GLUE \\
NLI train & 10,000 & 9,952 & 1--5 & SNLI/MultiNLI \\
extra\_train & 22,471 & 21,943 & 0--5 & STS12--16 \\
SICK-R & 9,927 & 9,927 & 1--5 & SICK-R \\
\bottomrule
\end{tabular}
\caption{STS training data composition.}
\end{table}

\textbf{Deduplication rule} (strictest, single-sentence level): any training
pair containing a sentence that appears in the STS-B dev/test sentence set is
removed, preventing the model from memorizing evaluation sentence
representations.

\textbf{AnglE Loss} ($\tau=0.50$, optimal for the albatross path, much higher
than the Rust path's 0.1):
\begin{equation}
\mathcal{L} = -\frac{1}{n} \sum_i \left[ y_i \log
\sigma\!\left(\frac{\cos\theta_i}{\tau}\right) + (1-y_i) \log\left(1 -
\sigma\!\left(\frac{\cos\theta_i}{\tau}\right)\right) \right],
\end{equation}
where $\cos\theta_i = \text{emb}_1 \cdot \text{emb}_2$ (embeddings are
L2-normalized), and $y_i$ is the min--max normalization of the score:
$y_i = (\text{score}_i - \min) / (\max - \min)$, unifying heterogeneous score
ranges (0--5, 1--5) from different datasets into $[0,1]$.

\textbf{5-seed ensemble}: train 5 models with different seeds (42, 123, 456,
789, 1024), average embeddings, and re-apply L2 normalization.

\subsection{Task 2: Topic Clustering}

\textbf{Data Source}:
\texttt{sklearn.fetch\_20newsgroups(}%
\allowbreak\texttt{subset='all', remove=('headers','footers','quotes'))},
removing headers/footers/quotes to prevent the model from memorizing
meta-information. Original 18,846 samples $\to$ filter 515 empty documents
$\to$ text-normalized hash deduplication removes 78 (0.43\%, cross-group
cross-posts) $\to$ 18,253 samples.

\textbf{Stratified split (70/15/15)}: train 12,777 / dev 2,738 / test 2,738.
Dev is used for model selection (best\_state), test for final evaluation
(held-out, not involved in training).

\textbf{Supervised contrastive learning}: pair construction 50\% same class
(score=1), 50\% different class (score=0); resample 80,000 pairs per epoch;
AnglE loss, $\tau=0.30$; 5-seed ensemble.

\textbf{Evaluation}: on the held-out test set, use projection + KMeans to
evaluate $v$-measure. \emph{Note}: this method uses class labels to train
the projector, making it supervised clustering---not directly comparable to
unsupervised MTEB benchmarks (e.g., bge-large's 0.57). The goal is to verify
the semantic potential of the hidden state, not to propose a new
unsupervised clustering method.

\subsection{Task 3: Task Classification}
\label{sec:classification-method}

\textbf{Method}: hidden state + MLP classifier (Top-K head selection + PCA
is included as a comparison baseline).
\begin{itemize}
  \item Hidden + MLP: $1024 \to 256 \to 4$, cross-entropy loss.
  \item Top-K head (comparison): evaluate each head's accuracy, select the
  top-8 heads' state (32,768 dims) $\to$ PCA to 256 dims $\to$ MLP.
\end{itemize}

\textbf{Strict split (70/15/15)}: train 70\% / dev 15\% (head selection +
early stopping) / test 15\% (final evaluation, held-out). Head selection is
ranked only on dev, never touching test, avoiding overfitting.

\section{Experiments}

\subsection{Experimental Setup}

\begin{itemize}
  \item \textbf{Backbone}: RWKV-7-g1 0.4B
  (\texttt{rwkv7-g1d-0.4b-20260210-ctx8192.pth}, BlinkDL/rwkv7-g1).
  \item \textbf{Inference engine}: albatross (official PyTorch + CUDA WKV
  kernel).
  \item \textbf{Training hardware}: CPU (Intel i7-12700K) + RTX 2080 Ti
  (backbone inference only).
  \item \textbf{Training framework}: PyTorch 2.0+ (CPU mode for projection
  training).
  \item \textbf{State layer}: L12 (experimentally determined optimal).
  \item \textbf{Max sequence length}: 128 tokens (truncated).
\end{itemize}

\subsection{Datasets}

\begin{table}[h]
\centering
\begin{tabular}{llll}
\toprule
Task & Dataset & Samples & Usage \\
\midrule
Clustering & 20NG (sklearn full-text, dedup) & 18,253 & train/dev/test (70/15/15) \\
Clustering (unsup.) & 20NG (MTEB short-text) & 59,545 & unsupervised evaluation \\
STS & STS-B train + NLI + extra + SICK-R & 46,898 pairs & training \\
STS & STS-B dev/test & 1,500 / 1,379 & evaluation \\
Classification & golden\_balanced & 16,751 & train/dev/test (70/15/15) \\
\bottomrule
\end{tabular}
\caption{Datasets used in the three tasks.}
\end{table}

\subsection{Main Results}

\subsubsection{STS Task}

\begin{table}[h]
\centering
\begin{tabular}{llr}
\toprule
Method & Training Data & Test Spearman \\
\midrule
Unsupervised hidden cosine & --- & 0.4600 \\
MLP + STS-B train (5.7k, not deduplicated) & 5,749 & 0.5818 \\
MLP + all data (46.9k, strictly deduplicated, baseline) & 46,898 & 0.8053 \\
\textbf{MLP + all data (optimal config)} & \textbf{46,898} &
\textbf{0.8188} \\
\bottomrule
\end{tabular}
\caption{STS results. Optimal config: hidden\_dim=1024, output\_dim=512,
dropout=0.1, $\tau=0.5$, 50 epochs, 5 seeds.}
\end{table}

\begin{table}[h]
\centering
\begin{tabular}{lrrr}
\toprule
Seed & dev & test & ensemble \\
\midrule
42 & 0.8125 & 0.7786 & 0.7786 \\
123 & 0.8071 & 0.7791 & 0.8006 \\
456 & 0.8021 & 0.7900 & 0.8127 \\
789 & 0.8022 & 0.7808 & 0.8157 \\
1024 & 0.8102 & 0.7874 & 0.8188 \\
\bottomrule
\end{tabular}
\caption{STS per-seed results.}
\end{table}

\textbf{Conclusion}: after strict deduplication (removing 1,249 pairs
overlapping with STS-B dev/test), baseline Spearman dropped from 0.8504 to
0.8053 ($-0.0451$), and the optimal config dropped from 0.8680 to 0.8188
($-0.0492$), confirming the impact of data leakage. The optimal config still
improves $+0.0135$ over the baseline, and the result is more honest, still
significantly outperforming the unsupervised baseline (0.46).

\textbf{Comparison with supervised embedding models}:

\begin{table}[h]
\centering
\begin{tabular}{llrl}
\toprule
Method & Type & STS-B Spearman & Parameters \\
\midrule
\textbf{Ours (RWKV-7 0.4B + projection, dedup)} & supervised &
\textbf{0.8188} & \textbf{3.15M (projector)} \\
bge-large-en-v1.5 \cite{bge} & supervised & $\sim$0.85 & 335M \\
all-MiniLM-L6-v2 \cite{minilm} & supervised & $\sim$0.86 & 22M \\
Unsupervised hidden cosine & unsupervised & 0.46 & --- \\
\bottomrule
\end{tabular}
\caption{Comparison with supervised embedding models. Supervised model data
from the MTEB leaderboard \cite{mteb}; training data and evaluation splits
may differ slightly across models, so this is an approximate comparison.}
\end{table}

Our method achieves STS-B Spearman of 0.8188 with only 3.15M projector
parameters (the 0.4B language model is frozen; only the projector is
trained), lower than bge-large-en-v1.5 (335M full fine-tuning) at
$\sim$0.85 and all-MiniLM-L6-v2 (22M) at $\sim$0.86. The gap is mainly
attributable to the representation capacity ceiling of the 0.4B language
model and the training data scale.

\subsubsection{Clustering Task (Supervised Projection + KMeans)}

This paper evaluates clustering on two complementary datasets: the
\textbf{MTEB short-text version} (59,545 samples, title-level; comparable
with the MTEB leaderboard, data download is difficult) and the
\textbf{sklearn full-text version} (18,253 samples, strict dedup +
stratified 70/15/15 split; reproducible by anyone, held-out evaluation).

\begin{table}[h]
\centering
\begin{tabular}{lrrr}
\toprule
Dataset & Samples & Unsupervised baseline & Supervised (SupCon) \\
\midrule
MTEB short-text & 59,545 & 0.2912 & \textbf{0.9506} (AnglE) \\
sklearn full-text & 18,253 & 0.4434 $\pm$ 0.0146 & \textbf{0.6660} \\
\bottomrule
\end{tabular}
\caption{Clustering main results on two datasets.}
\end{table}

\begin{table}[h]
\centering
\small
\begin{tabular}{llr}
\toprule
Method & Training type & Test $v$-measure \\
\midrule
Hidden + standardize + KMeans (baseline, 10 seeds) & unsupervised &
0.4434 $\pm$ 0.0146 \\
STS projection transfer (failed) & supervised (STS) & 0.1424 \\
Supervised contrastive (AnglE, $\tau$=0.3, 80k pairs) & supervised
(clustering) & 0.6373 \\
\textbf{Supervised contrastive (SupCon, $\tau$=0.07, batch=320)} &
\textbf{supervised (clustering)} & \textbf{0.6660} \\
\bottomrule
\end{tabular}
\caption{sklearn full-text detailed results.}
\end{table}

\begin{table}[h]
\centering
\begin{tabular}{lrrr}
\toprule
Seed & dev $v$ & test $v$ & ensemble \\
\midrule
42 & 0.6351 & 0.6402 & 0.6402 \\
123 & 0.6374 & 0.6378 & 0.6528 \\
456 & 0.6281 & 0.6458 & 0.6614 \\
789 & 0.6338 & 0.6473 & 0.6647 \\
1024 & 0.6355 & 0.6418 & \textbf{0.6660} \\
\bottomrule
\end{tabular}
\caption{sklearn full-text per-seed results (SupCon).}
\end{table}

\textbf{Configuration}: SupCon loss, $\tau=0.07$, batch=320 (PK sampler, 20
classes $\times$ 16), dropout=0.1, 30 epochs, 5-seed ensemble.

\textbf{Why the MTEB short-text $v$-measure $=0.9506$ has no data leakage}.
Since this result is notably higher than the sklearn full-text version
(0.6660), we detail the data split and evaluation pipeline here to confirm
no leakage. A common confusion must be clarified first: an earlier version
mistakenly labeled ``Projection + standardize + KMeans $=$ 0.8466'' (a
supervised-projection result) as the unsupervised baseline. After rerunning
the unsupervised baseline (10-seed KMeans) and the supervised projection
(5 seeds), the \textbf{true unsupervised baseline of the MTEB short-text
version is 0.2912}. The 0.8466 is a supervised-projection result, not an
unsupervised baseline.

\begin{enumerate}
  \item \textbf{Data split (64/16/20 stratified)}: the 59,545 samples are
  split via stratified sampling by class labels into train (38,109; only for
  projector training), dev (9,527; only for best\_state model selection),
  and test (11,909; fully held-out, never used in training or model
  selection).
  \item \textbf{Evaluation pipeline (unsupervised clustering + supervised
  metric)}: the projector is trained only on train; best\_state is selected
  by $v$-measure on dev. Final evaluation runs on test: project test hidden
  states to 128-d via the trained projector, then cluster with
  \textbf{KMeans (unsupervised, $k{=}20$)}. \textbf{Class labels are used
  only to compute the $v$-measure metric, not in the KMeans clustering
  step}, so there is no label leakage.
  \item \textbf{Why the high result is reasonable}: the MTEB short-text
  unsupervised baseline is only 0.2912, indicating severe hidden-state
  anisotropy. Supervised projection lifts $v$-measure from 0.2912 to 0.9506
  ($+217\%$), an even larger gain than the $+50\%$ on the sklearn full-text
  version. Short-text samples have fewer tokens and more concentrated
  semantics, making it easier for supervised projection to learn clear
  inter-class separation boundaries.
  \item \textbf{No-overfitting evidence}: dev $v=0.9497$ is close to test
  $v=0.9506$ (after the 5-seed ensemble), indicating the best\_state
  selected on dev performs equally well on test.
\end{enumerate}

\textbf{Analysis of differences between datasets}: (1) the sklearn
full-text supervised result (0.6660) is much lower than MTEB (0.9506), due
to less training data (12k vs 38k) plus harder long-text clustering; (2) but
the sklearn full-text unsupervised baseline (0.4434) is higher than MTEB
(0.2912), because removing headers/footers/quotes yields cleaner samples;
(3) both datasets validate the core thesis: supervised projection
significantly outperforms unsupervised methods (sklearn $+50\%$, MTEB
$+217\%$).

\subsubsection{Unsupervised Clustering Methods Comparison}
\label{sec:unsupervised-comparison}

To validate the claim that unsupervised methods cannot extract clustering
structure from the hidden state, we systematically compared 7 categories and
30+ unsupervised methods on the official MTEB 20NG short-text version
(59,545 samples, title-level):

\begin{table}[h]
\centering
\begin{tabular}{llr}
\toprule
Method category & Best method & Best $v$-measure \\
\midrule
Hidden + KMeans (baseline) & L23 + standardize & 0.2912 \\
Nonlinear feature transform & L23 + sign$\cdot$sqrt + standardize & 0.3181 \\
Nonlinear dim.\ reduction (UMAP) & UMAP(32) + KMeans & 0.3263 \\
Multi-layer concatenation & Last3 (L16+L20+L23) & 0.2553 \\
Layer difference & L23$-$L0 + sign$\cdot$sqrt & 0.3087 \\
Whitening / top-$k$ PC removal & PCA whitening 256 & 0.1629 \\
WKV state aggregation stats & row\_mean + col\_mean & 0.1020 \\
DeepCluster self-supervised iter. & Round 1 ensemble & 0.3021 \\
\bottomrule
\end{tabular}
\caption{Unsupervised methods comparison on the MTEB short-text version.}
\end{table}

\textbf{Analysis}: (1) linear methods cap at $\sim$0.32, far below the
unsupervised MTEB SOTA of 0.57; (2) nonlinear dimensionality reduction
(UMAP 0.3263) offers no breakthrough, indicating the problem is not in the
reduction method; (3) multi-layer concatenation (0.26) is worse than
single-layer L23 (0.32)---shallow-layer noise dilutes deep-layer semantics;
(4) DeepCluster pseudo-labels cannot self-improve (0.30$\to$0.30); (5) WKV
state aggregation statistics peak at 0.10---the WKV state's value range is
too small (std $= 0.13$); clustering information resides in the hidden
state, not the WKV state.

\textbf{Root cause}: the 0.4B language model's optimization objective is
next-token prediction, not clustering. The semantic information in the
hidden state requires nonlinear transformation (supervised projection MLP)
to unlock. Unsupervised methods cap at 0.33 (MTEB short-text) / 0.44
(sklearn full-text) vs supervised projection 0.6660 (sklearn full-text,
SupCon) / 0.9506 (MTEB, AnglE).

\subsubsection{Classification Task}

\begin{table}[h]
\centering
\begin{tabular}{lrr}
\toprule
Method & dev acc & test acc \\
\midrule
\textbf{Hidden + MLP} & 0.9381 & \textbf{0.9325} \\
Top-8 head + PCA256 + MLP & 0.9247 & 0.9208 \\
\bottomrule
\end{tabular}
\caption{Classification results.}
\end{table}

\textbf{Analysis}: hidden + MLP outperforms top-$k$ head + PCA (0.9325 $>$
0.9208), indicating that the albatross hidden state already possesses good
task separability---directly training a classifier on hidden is sufficient;
top-$k$ head selection + PCA dimensionality reduction actually loses some
information. dev and test accuracies are close (0.9381 vs 0.9325),
indicating no overfitting. This is consistent with the core insight in
\S\ref{sec:insight}---the hidden state contains semantic information that
can be extracted with a simple MLP (classification is a discriminative task
that does not require projection to a semantic space).

\subsection{Ablation Studies}
\label{sec:ablation}

\subsubsection{STS Projection Transfer to Clustering (Failed)}

This experiment uses a sampled subset of 2,000 samples (100 per class),
different from the full test set (11,909 samples):

\begin{table}[h]
\centering
\begin{tabular}{lr}
\toprule
Method & $v$-measure \\
\midrule
Baseline (hidden + standardize, sampled 2,000) & 0.3426 \\
STS projection + KMeans & 0.1424 \\
STS projection + standardize + KMeans & 0.1305 \\
STS projection + PCA(64) + KMeans & 0.1284 \\
Hidden + STS projection concat & 0.2756 \\
\bottomrule
\end{tabular}
\caption{STS projector transfer to clustering fails.}
\end{table}

\textbf{Failure reason}: STS learns ``similarity ranking'' (relative
distance); after L2 normalization, embedding std $= 0.0884$ is too small,
compressing inter-class distances. Clustering requires ``class separation''
(absolute distance)---the two objectives differ.

\subsubsection{Training Data Scale Impact (STS)}

\begin{table}[h]
\centering
\begin{tabular}{lrr}
\toprule
Training data & pairs & Test Spearman \\
\midrule
STS-B train only & 5,749 & 0.5818 \\
+ NLI + extra + SICK-R (not deduplicated) & 48,147 & 0.8504 (leaked) \\
+ NLI + extra + SICK-R (strictly deduplicated, baseline) & 46,898 & 0.8053 \\
+ NLI + extra + SICK-R (strictly deduplicated, optimal) & 46,898 &
\textbf{0.8188} \\
\bottomrule
\end{tabular}
\caption{Training data scale impact on STS.}
\end{table}

\textbf{Conclusion}: data volume is the key bottleneck for the STS task; an
8$\times$ increase brings a $+40\%$ improvement (0.5818$\to$0.8188). The
non-deduplicated version 0.8504 was inflated due to training-data sentence
overlap with dev/test; 0.8188 (optimal config) is the true generalization
result.

\subsection{Summary of Failed Directions}

\begin{table}[h]
\centering
\begin{tabular}{lll}
\toprule
Method & Result & Failure reason \\
\midrule
STS projection transfer to clustering & 0.14 & ranking vs separation \\
Unsupervised KMeans & 0.29 & severe anisotropy \\
Unsupervised hidden cosine (STS) & 0.46 & anisotropy \\
UMAP nonlinear dim.\ reduction & 0.33 & problem not in reduction \\
DeepCluster self-supervised iter. & 0.30 & pseudo-labels too weak \\
Multi-layer hidden concatenation & 0.26 & shallow noise dilutes \\
Whitening / top-$k$ PC removal & 0.16 & degrades performance \\
Pure WKV state (Q-Readout) & 0.11 & state range small (std 0.13) \\
WKV state aggregation stats & 0.10 & no clustering info \\
Disable v\_first mechanism & 0.16 & state std only 0.16 \\
fp32 state (vs fp16) & 0.13 & precision not the root cause \\
\bottomrule
\end{tabular}
\caption{Summary of failed directions.}
\end{table}

\section{Production Deployment: Agentic Model Routing}
\label{sec:production}

The classification pipeline of \S\ref{sec:classification-method} ships as
the smart router of \textbf{ai00-x}, an open-source AI assistant desktop
client with a resident local RWKV engine. This section reports the design,
two deployment findings we believe are independently useful
(\S\ref{sec:logits-fail}, \S\ref{sec:language-shortcut}), and production
measurements on a 2.9B backbone.

\subsection{System Design}

Each conversational turn is routed across four complexity tiers, following
the agentic-routing philosophy of SquillaRouter \cite{squillarouter}:

\begin{table}[h]
\centering
\begin{tabular}{lll}
\toprule
Tier & Semantics & Default target \\
\midrule
R0 & greetings, acknowledgments & local RWKV (zero API cost) \\
R1 & simple single-step tasks & local RWKV \\
R2 & code generation, multi-step reasoning & mid-tier model (user-configured) \\
R3 & debugging, long-context, high-risk & primary model (user-configured) \\
\bottomrule
\end{tabular}
\caption{Four-tier routing semantics.}
\end{table}

The pipeline: a zero-cost rule short-circuit intercepts trivial
acknowledgments; otherwise the request is classified by RWKV; a
post-processing rule stack applies a safety upgrade (predicted R0/R1 but
$P(\text{R2}) + P(\text{R3}) > 0.45 \Rightarrow \text{R2}$) and a
sticky-tier rule within a session; the tier maps to a model; any failure
(engine not ready, head missing, dimension mismatch, timeout) falls back to
the primary model.

The classifier runs \textbf{inside the already-resident inference engine}:
tokenize (128-token truncation), prefill from a zero state on a dedicated
slot isolated from the 16 generation slots, mean-pool the last-layer hidden
state, and score it with the trained head. Unlike SquillaRouter
\cite{squillarouter}---390 hand-crafted features, a LightGBM/ONNX runtime,
and a separate on-device embedding model---this design adds \textbf{zero
runtime dependencies}: the head is a 0.66M-parameter matrix multiply (0.32
ms on CPU) over features of the model that must stay resident anyway for
local generation.

\subsection{Next-Token Logits Fail; State Embeddings Work}
\label{sec:logits-fail}

Before adopting state embeddings, we tested the seemingly simpler
alternative: prompt the base model to answer with a tier label and read
next-token logits. On the untuned 2.9B base model this fails
categorically---a raw-prompt variant collapses to R0 (30\% accuracy), a
User:/Assistant:-wrapped variant collapses to R2 (28\%). The next-token
distribution of a small base model carries no trustworthy classification
signal, while its hidden state yields 96\%+
(\S\ref{sec:results-29b}). This mirrors the paper's core claim at the
output layer: the information is in the state, not on the surface.

\subsection{A Language-Shortcut Pitfall}
\label{sec:language-shortcut}

The golden\_balanced training set contains 1,783 Chinese samples---all
labeled R0. The first trained head duly learned a ``Chinese $\to$ R0''
shortcut: Chinese R1/R2/R3 requests collapsed to R0 with probability 1.00,
while English traffic still scored $\sim$96\%. Purely English evaluation
cannot detect this failure. We generated 857 Chinese R1/R2/R3
template-composed augmentation samples and merged them into the training
split only (held-out sets untouched). After retraining, a hand-crafted
Chinese--English mixed set rose from 45\% to 100\%, and the English-dominant
held-out set from 96.2\% to 97.1\%---no regression. We flag this as a
general data-audit item for tier classifiers trained on mixed-language
traffic.

\subsection{Results at 2.9B}
\label{sec:results-29b}

Scaling the frozen backbone from 0.4B (hidden 1024) to a 2.9B int8-quantized
model (hidden 2560), same head architecture and training protocol (70/15/15
stratified split; AdamW 1e-3, cosine schedule, 30 epochs, early stopping on
dev):

\begin{table}[h]
\centering
\begin{tabular}{lrr}
\toprule
Backbone & dev acc & test acc \\
\midrule
0.4B fp16 (\S4.3.4) & 0.9381 & 0.9325 \\
2.9B int8 + zh augmentation & 0.969 & \textbf{0.960} \\
\bottomrule
\end{tabular}
\caption{Classification accuracy across backbones.}
\end{table}

The comparison jointly changes backbone scale and augmentation; the
unaugmented 2.9B head already classifies English traffic at $\sim$96\% but
collapses on Chinese R1--R3 (\S\ref{sec:language-shortcut}), so most of the
gain is attributable to the backbone, with augmentation repairing the
Chinese failure mode.

End-to-end through the deployed engine: 97.1\% on an independent 314-sample
held-out set (per-class: R0 100\% / R1 95\% / R2 96\% / R3 100\%) and 100\%
on a 40-sample hand-crafted Chinese--English mixed set. The Python-trained
weights and the Rust inference port agree exactly (max probability
difference 0.000000).

\textbf{Latency} (RTX 2080 Ti; full pipeline tokenize $\to$ zero-state
prefill $\to$ mean-pool $\to$ head; 250 samples after warm-up): mean 77 ms,
p50 79 ms, p95 81 ms; throughput 12.4 classifications/s. Latency is nearly
insensitive to request length (16$\to$128 tokens: $+12\%$) because the
sequence-path GEMM pads the token dimension to 256. Stage breakdown at 128
tokens: prefill 93.8\%, zero-state restore 5.8\%, head 0.4\%. Since R0/R1
turns are answered by the same local model, $\sim$80 ms is negligible
against local generation (seconds) or a cloud first token (500 ms+), and the
3 s routing timeout is never approached.

\subsection{Deployment Notes}

\begin{itemize}
  \item \textbf{Graceful degradation}: a missing, corrupted, or
  dimension-mismatched head (e.g., after the user swaps the local model)
  silently disables classification and traffic falls back to the primary
  model; the system runs with or without the router.
  \item \textbf{Distribution}: the 14.2 MB head ships through the same
  content-addressed model manifest as the backbone (sha256-verified,
  mirrored to ModelScope/HuggingFace), adding $\sim$14 MB to first-launch
  downloads.
  \item \textbf{Engine prerequisites}: two inference-engine bugs had to be
  fixed before the latency above was reachable---a kernel-cache invalidation
  that recompiled CUDA kernels ($\sim$0.5 s) on every new sequence length,
  and a VRAM leak on variable-length prefills. Both affect any
  variable-length prefill workload, not only classification.
\end{itemize}

\subsection{Comparison with SquillaRouter}

\begin{table}[h]
\centering
\small
\begin{tabular}{lll}
\toprule
 & SquillaRouter \cite{squillarouter} & Ours \\
\midrule
Features & 390 hand-crafted + external embeddings & native hidden state (2560-d) \\
Classifier & LightGBM + ONNX & 0.66M-param MLP, pure Rust \\
Extra runtime deps & LightGBM, ONNX, embedding model & none \\
Headline result & cost $-90\%$ / $-43\%$ on agentic benchmarks & test acc 0.960,
97.1\% end-to-end, 77 ms \\
\bottomrule
\end{tabular}
\caption{Comparison with SquillaRouter.}
\end{table}

The two results are complementary rather than directly comparable:
SquillaRouter reports cost--quality frontiers on agentic benchmarks; we
report tier-classification accuracy and latency. The shared element is the
four-tier philosophy; the engineering difference is that our router needs
nothing beyond the local model the assistant already runs, with R0/R1
served at zero API cost. SquillaRouter's self-improving data flywheel---
online capture of routing outcomes and periodic router retraining---remains
future work for our deployment.

\section{Discussion}

\subsection{Core Insight: Supervised Projection Unlocks Hidden Semantic
Potential}
\label{sec:insight}

The anisotropy of the albatross hidden state (unsupervised STS only 0.46)
is not a feature defect, but evidence that \textbf{unsupervised methods
cannot extract nonlinear structure}. Supervised MLP classification achieving
0.94 proves the features contain information; only a supervised projector
(MLP + AnglE loss) is needed to map it to a linearly separable semantic
space.

\subsection{Task-Specific Projectors Cannot Be Mixed}

\begin{table}[h]
\centering
\begin{tabular}{lll}
\toprule
Task & Objective & Projector behavior \\
\midrule
STS & similarity ranking & preserves relative distances \\
Clustering & class separation & amplifies inter-class distances \\
\bottomrule
\end{tabular}
\caption{Task-specific projector behaviors.}
\end{table}

STS projector transfer to clustering fails (0.14 $<$ 0.34 baseline),
proving that \textbf{projectors must align with task objectives}.

\subsection{Data Scale Is Key}

STS task expansion from 5.7k to 46.9k (8$\times$) brings a $+39\%$
improvement, indicating that the albatross hidden state's semantic
information requires sufficient data to extract through supervised
learning. This is consistent with SimCSE's finding that contrastive
learning requires large amounts of samples \cite{simcse}.

\subsection{Albatross vs Rust (web-rwkv)}

The optimal $\tau=0.50$ for the albatross path is much higher than the Rust
path's 0.1, rooted in hidden value-range differences (albatross std $=1.70$
vs Rust std $=3.54$). This work abandons comparison with Rust, using the
albatross official implementation as the standard to avoid modifying the
inference engine.

\section{Conclusion}

This paper proposes a supervised projection-based framework for RWKV-7
hidden state semantic embedding extraction, adopting a strict experimental
paradigm (STS training data deduplication; clustering on the sklearn
full-text version with strict deduplication + stratified split;
classification with an independent test set), achieving the following
results on three standard tasks:

\begin{itemize}
  \item \textbf{Semantic similarity (supervised)}: Spearman $=0.8188$ (after
  deduplication), lower than bge-large-en-v1.5 ($\sim$0.85, 335M) and
  all-MiniLM-L6-v2 ($\sim$0.86, 22M), but with only 3.15M projector
  parameters. The gap is mainly attributable to the representation capacity
  ceiling of the 0.4B language model.
  \item \textbf{Topic clustering (supervised)}: validated on two
  datasets---MTEB short-text $v$-measure $=0.9506$ ($+217\%$ over the
  unsupervised baseline 0.2912); sklearn full-text SupCon loss
  $v$-measure $=0.6660$ ($+50\%$ over the unsupervised baseline 0.4434).
  SupCon improves $+0.0287$ over AnglE on the sklearn version.
  \item \textbf{Task classification}: test acc $=0.9325$ (independent test
  set, head selection uses only dev).
  \item \textbf{Production routing}
  (\S\ref{sec:production}): the same classification head, scaled to a 2.9B
  int8 backbone with Chinese augmentation and hardened by a rule stack,
  routes live traffic in a shipped desktop assistant---0.960 test accuracy,
  97.1\% end-to-end, 77 ms mean latency, zero extra runtime dependencies,
  R0/R1 served locally at zero API cost.
\end{itemize}

Core insights: the albatross hidden state contains semantic information but
requires supervised projection to unlock; task-specific projectors cannot
be mixed; data scale is the key bottleneck. The WKV state shows much lower
clustering performance than the hidden state (0.11 vs 0.44), indicating the
hidden state is more suitable for semantic extraction. All offline methods
are based on a 0.4B model, the albatross inference engine (source code
unmodified), CPU trainable with $\sim$3.15M parameters, suitable for edge
deployment; the production router (\S\ref{sec:production}) additionally
validates the classification pipeline on a 2.9B int8 backbone in live
traffic.

\paragraph{Honest disclosure.} The early-version reported STS 0.8504 was
inflated due to data leakage (STS training data overlapped with STS-B test
by 1,249 pairs); the clustering 0.8466 was actually a ``projection +
standardize + KMeans'' result that was mistakenly labeled as the
unsupervised baseline. After rerunning the unsupervised baseline (10 seeds)
and the supervised projection (5 seeds), and rerunning all three supervised
clustering experiments for verification, the true unsupervised baseline of
the MTEB short-text version is 0.2912. This version has been strictly
deduplicated, uses held-out evaluation, and corrects the baseline-labeling
error, giving more honest results.

\begin{thebibliography}{15}

\bibitem{rwkv7} RWKV. RWKV-7: Linear RNN with Delta Rule.
\url{https://www.rwkv.cn/}.

\bibitem{deltarule} Delta Rule. Wikipedia.
\url{https://en.wikipedia.org/wiki/Delta_rule}.

\bibitem{albatross} BlinkDL. Albatross: Efficient RWKV-7 Inference Engine.
\url{https://github.com/BlinkDL/Albatross}.

\bibitem{mteb} Muennighoff, N. et al. MTEB: Massive Text Embedding
Benchmark. arXiv:2210.07316.

\bibitem{stsb} Cer, D. et al. SemEval-2017 Task 1: Semantic Textual
Similarity Multilingual and Crosslingual Focused Evaluation. SemEval 2017.

\bibitem{20ng} 20 Newsgroups.
\url{http://qwone.com/~jason/20Newsgroups/}.

\bibitem{sick} Marelli, M. et al. SICK: A Sentiment Inference Corpus.
LREC 2014.

\bibitem{angle} Li, X. et al. AnglE-optimized Text Embeddings.
arXiv:2309.12871.

\bibitem{simcse} Gao, T., Yao, X., and Chen, D. SimCSE: Simple Contrastive
Learning of Sentence Embeddings. EMNLP 2021.

\bibitem{whitening} Su, J. et al. Whitening Sentence Representations for
Better Semantics and Faster Retrieval. arXiv:2103.15316.

\bibitem{sbert} Reimers, N. and Gurevych, I. Sentence-BERT: Sentence
Embeddings using Siamese BERT-Networks. EMNLP 2019.

\bibitem{bge} Xiao, S. et al. BAAI/bge-large-en-v1.5.
\url{https://huggingface.co/BAAI/bge-large-en-v1.5}.

\bibitem{minilm} Wang, W. et al.
sentence-transformers/all-MiniLM-L6-v2.
\url{https://huggingface.co/sentence-transformers/all-MiniLM-L6-v2}.

\bibitem{mtebleaderboard} MTEB Leaderboard.
\url{https://huggingface.co/spaces/mteb/leaderboard}.

\bibitem{squillarouter} Liu, X. et al. Agentic Routing: The Harness-Native
Data Flywheel. arXiv:2607.11399, 2026.

\end{thebibliography}

\appendix

\section{Reproduction Guide}

\paragraph{Requirements.} Python 3.10+ with \texttt{uv}; a CUDA GPU (for
albatross inference); on Windows, MSVC is required for CUDA extension
compilation.

\paragraph{Reproduction commands.}

\begin{verbatim}
# 1. Download model and albatross source
uv run --project ../../scripts python 00_setup.py

# 2. Extract features
./run_with_msvc.bat extract_features.py --task sts --sts-subdir
    sts_dedup --batch-size 16 --max-length 128
./run_with_msvc.bat extract_features.py --task cluster_sklearn
    --batch-size 16 --max-length 128
./run_with_msvc.bat extract_features.py --task classification
    --batch-size 16 --max-length 128 --limit 8000

# 3. Run three tasks
uv run python 02_sts_similarity.py --data-dir ../../data/sts_dedup
    --device cuda --n-epochs 50 --hidden-dim 1024
    --output-dim 512 --dropout 0.1        # STS: 0.8188
uv run python 07_cluster_supcon.py --device cuda   # Clustering: 0.6660
uv run python 03_classification.py          # Classification: 0.9325
\end{verbatim}

\paragraph{Expected output.}
\begin{verbatim}
STS:          5seed ensemble: Spearman = 0.8188
Clustering:   5seed ensemble: v_measure = 0.6660 (sklearn SupCon)
                              / 0.9506 (MTEB AnglE)
Classification: test_acc = 0.9325
\end{verbatim}

\paragraph{Production router (\S\ref{sec:production}).}
Feature extraction for the 2.9B backbone runs through the deployed Rust
engine (\texttt{extract\_router\_features.rs} in the rwkv-rsv repository);
training and Chinese augmentation live in the ai00-x client repository
(\texttt{scripts/router\_head/}: \texttt{train\_mlp.py},
\texttt{gen\_zh\_augment.py}), using the same protocol as above. Production
measurements (77 ms mean latency, 12.4 classifications/s) are from the
deployed desktop client on an RTX 2080 Ti.

\end{document}
